\documentclass[10pt]{article}
\usepackage[margin=20mm]{geometry}
\usepackage{amsmath,amssymb,fontspec,unicode-math,xurl,hyperref}
\IfFontExistsTF{DejaVu Serif}{\setmainfont{DejaVu Serif}}{\setmainfont{Latin Modern Roman}}
\setmathfont{latinmodern-math.otf}
\setlength{\parindent}{0pt}\setlength{\parskip}{6pt}\sloppy
\begin{document}
Abstract: This review checks the published representation-based weak-to-strong generalization analysis, its Gaussian benign-overfitting example, and the scope of the empirical predictor. A Gaussian family satisfying the stated example premises fails the kernel-isotropy condition needed by the supplied proof. Its error is of order \(\gamma^2\), whereas the definition requires \(o(\gamma^2)\). This establishes a Major defect in the applicability derivation underlying the advertised benign-overfitting application, without refuting the main theorem under its actual assumptions or the example’s final risk formulas. A separate finite perturbation inequality is false, although a verified denominator correction preserves the asymptotic uses inspected here. No author code or original experiments were executed. Scientific adjudication remains pending.\par
Author: AI-assisted technical review\par
Date:\par
Title: Technical Verification of
Representations Shape Weak-to-Strong Generalization\par
\# Identity, reviewed version, and scope\par
The reviewed paper is Yihao Xue, Jiping Li, and Baharan Mirzasoleiman, Representations Shape Weak-to-Strong Generalization: Theoretical Insights and Empirical Predictions, ICML 2025, Proceedings of Machine Learning Research 267, pp. 70120–70159, publisher version dated 6 October 2025; internal identifiers are 0142 and I25-43511, and OpenReview identifier is `ypEW077kle` [Retained latex: \textbackslash{}cite\{paper\}]. All page references below are one-based pages of the 40-page PDF; its printed local page numbers agree with these PDF numbers. The staged original and source manifest were read directly. The actual original PDF hash is recorded in Appendix 6 (\#sec:reproduction) and the verification records.\par
The principal inspected chain is Definition 3.3 \(\to\) Assumption 3.7 \(\to\) Theorem 3.8, including Appendix A’s ridge-regression factorization and Lemmas A.5–A.13. The review also inspected Theorem 4.1, Example 4.2 and Corollary 4.3 with Appendix B.3, the Gaussian verification in Appendix C.2.1, the construction in Theorem 3.6/Appendix C.3, and the definitions, training protocol and tuning discussion for the empirical metric in Section 5 and Appendix D.1–D.2.\par
Problem 1 disproves the claim that the stated Gaussian example meets the prescribed decomposability conditions over its full parameter range. It does not supply a counterexample satisfying Assumption 3.7 to Theorem 3.8. The exact high-probability argument, rather than finite simulation, establishes the defect.\par
The provided manifest records no located paper-specific implementation and an OpenReview attachment access challenge. These are inherited source-availability facts, not independently resolved by new network requests. Experiments, pretrained models, raw predictions, and numerical figure correlations were not reproduced. No later paper revision or erratum was supplied or checked; first-discovery status is not asserted.\par
\# Definitions and inspected conditional reasoning\par
Equation reference: eq:required. For a representation \(h\), the paper defines \(\Sigma(h)=\mathbb E[h(x)h(x)^\top]\) and an unnormalized sample Gram matrix \(\widehat K(h)\). Definition 3.3 decomposes the representation space into an eigenspace \(V\) of \(\Sigma(h)\) and its orthogonal complement. With \[\gamma=\min(\widehat\gamma,\widetilde\gamma),
 \qquad \rho=\lambda_{\min,\ne0}(\Sigma(\Pi_Vh)),\] the kernel-isotropy condition on PDF p. 4 is \[
 \left\|\frac{\widehat K(\Pi_{V^\perp}h)}{\widehat n}
       -\widehat\gamma I\right\|_{\mathrm{op}}
   =o(\gamma^2+\delta^2),\] and similarly for the other split. The remaining conditions bound representation and label magnitudes, require concentration on \(V\), small cross-split nonprincipal Gram matrices, and vanishing nonprincipal population covariance.\par
Equation reference: eq:main. Theorem 3.8, PDF p. 5, assumes these conditions for both models and \(\delta_j\le\beta_j=O(1)\). With \(G_j=\widehat K(\Pi_{V_j}h_j)/\widehat n\), set \[P_w=G_w(G_w+(\beta_w+\widetilde\gamma_w)I)^{-1},
 \qquad P_s=G_s(G_s+(\beta_s+\widehat\gamma_s)I)^{-1}.\] The conclusion is \[
 \mathrm{PredGap}=\left\|P_s(I-P_w)\frac{\widehat y}{\sqrt{\widehat n}}\right\|^2\pm o(1).\] The ridge solutions and population quadratic form in Appendix A.4 give the appropriate starting expression. Lemma A.12 approximates weak predictions on the strong-training split; Lemma A.13 transfers the strong population quadratic form to its principal kernel. Subject to genuine decomposability and the corrected inverse bound in Problem 2, this inspected structure does not establish a new defect in [eq:main] (\#eq:main).\par
\# Substantive findings\par
Problem 1 (The Gaussian applicability proof misses the required isotropy rate).  The Gaussian premises in Appendix C.2.1 and Example 4.2 permit a family for which [eq:required] (\#eq:required) fails. The published route from that example to the advertised benign-overfitting guarantee therefore does not cover the full stated parameter range.\par
\#\# Original claim, location, and dependency\par
Equation reference: eq:paper34. Example 3.5, PDF pp. 4–5, and its explicit Gaussian specialization in Appendix C.2.1, PDF p. 29, allow \(\widetilde n=\Theta(\widehat n)\), \(d=\omega(\widehat n^2)\), \(\sigma^2=O(\widehat n)\) and \(\widehat n=\omega(k^2)\). The Gaussian covariance is \[\Sigma=\operatorname{diag}\left(I_k,\frac{\sigma^2}{d-k}I_{d-k}\right).\] Appendix C.2.1 states that these representations are \((0,\sigma^2/\widehat n,\sigma^2/\widetilde n)\)-decomposable. Its step 3, PDF p. 32, equation (34), correctly writes the nonprincipal Gram error as \[
 \frac{\sigma^2}{\widehat n}
 \left\|\frac{Z^\top Z}{d-k}-I\right\|_{\mathrm{op}}.\] The following assertion that the required conclusion follows by taking the concentration parameter \(\delta_2=\widehat n^{-1/4}\) does not verify the stronger \(o(\gamma^2)\) scale.\par
This claim is used in Example 4.2, PDF p. 6, and explicitly in Appendix B.3.2, PDF pp. 24–25: the Gaussian representations are declared decomposable, allowing Lemmas A.12, B.5 and Theorem 3.8 to determine the risks in Corollary 4.3. This dependency connects the defect to the core explanation of benign overfitting.\par
\#\# Admissible family and all example premises\par
Equation reference: eq:parameters. Let \(n\to\infty\) through integers, and take \[
 \widehat n=\widetilde n=n,\quad k=1,\quad
 m=d-k=n^3,\quad \sigma^2=1,\quad
 \eta=\frac35,\quad \beta_w=\beta_s=n^{-2}.\] For each independent sample, let \(y,z\sim N(0,1)\) independently, and take \[h_w=(\sqrt\eta\,y+\sqrt{1-\eta}\,z,\ \xi_w^\top)^\top,
 \qquad h_s=(y,\ \xi_s^\top)^\top,\] where \(\xi_w,\xi_s\sim N(0,I_m/m)\) are independent of each other and of \(y,z\). Samples and the two training splits are independent.\par
These choices satisfy every explicit premise of the Gaussian specialization and Example 4.2: first-coordinate variance one; nonprincipal-coordinate variance \(\sigma^2/(d-1)=1/m\); independent coordinates within each representation; \(\mathbb E[y^2]=1\); \(\eta\in(0,1)\); equal sample sizes tending to infinity; \(d/n^2\to\infty\); nonzero \(\sigma^2=o(n)\); and \(\beta_j=o(\sigma^2/n)\). The first coordinate is an exact population eigenspace with \(\rho=1\). We target the unambiguous Gaussian specialization rather than rely on the inconsistent variance-versus-standard-deviation convention for the sub-Gaussian parameter between Example 3.5 and Definition C.5.\par
The other decomposability conditions do not explain the failure. Principal empirical moments and empirical label second moments concentrate by finite Gaussian moments. The normalized sample covariance is bounded. The nonprincipal population covariance is \(1/m=o(1/n)\). For independent noise matrices, \(B=\widehat\Xi^\top\widetilde\Xi/n\) satisfies \[\mathbb E\|B\|_F^2=\frac1m,\qquad
 \mathbb P\{\|B\|_{\mathrm{op}}>\varepsilon/n\}
 \le \frac{n^2}{m\varepsilon^2}\longrightarrow0.\] Thus the cross-split condition holds, whereas the sharper within-split requirement [eq:required] (\#eq:required) does not.\par
\#\# Exact independent calculation\par
Equation reference: eq:exactcolumn, eq:lower. For either model, write \(\Xi=[\xi_1,\ldots,\xi_n]\in\mathbb R^{m\times n}\) and \[A_n=\frac{\Xi^\top\Xi-I_n}{n},\qquad \gamma=\frac1n.\] The required conclusion is \(\|A_n\|_{\mathrm{op}}/\gamma^2\to0\) in probability. Conditionally on \(\xi_1\), the off-diagonal entries in the first column are independent with \[(A_n)_{i1}=\xi_i^\top\xi_1/n
 \sim N\left(0,\frac{\|\xi_1\|^2}{mn^2}\right),\quad i=2,\ldots,n.\] Consequently, for independent \(U_m\sim\chi_m^2/m\) and \(V_{n-1}\sim\chi_{n-1}^2\), \[
 \sum_{i=2}^n(A_n)_{i1}^2\stackrel{d}{=}\frac{U_mV_{n-1}}{mn^2}.\] The operator norm is at least any column norm. Therefore its ratio to \(\gamma^2\) is bounded below by a random variable distributed as \[
 \frac{\left(\sum_{i=2}^n(A_n)_{i1}^2\right)^{1/2}}{\gamma^2}
 \stackrel{d}{=}n\sqrt{\frac{U_mV_{n-1}}m}
 =\sqrt{U_{n^3}\frac{V_{n-1}}n}
 \xrightarrow{\mathbb P}1.\] This uses only chi-squared means and variances. In particular, \(\mathbb P\{\|A_n\|_{\mathrm{op}}/\gamma^2>1/2\}\to1\), contradicting the asserted little-\(o\) rate.\par
Equation [eq:lower] (\#eq:lower) proves a lower bound, so the issue is not merely a loose printed upper bound. Independent finite diagnostics gave \(\|A_n\|_{\mathrm{op}}/\gamma^2=2.138,1.902,1.762\) for \(n=8,16,32\). These illustrate the analytic result and do not substitute for it.\par
\#\# Alternative readings and counterevidence\par
Replacing \(o(\gamma^2)\) by \(o(\gamma)\) avoids this contradiction, since \(\|A_n\|_{\mathrm{op}}/\gamma\) vanishes. The square is nevertheless explicit in the rendered source on PDF p. 4 and in Lemma A.5’s absolute inverse-error argument. Such a replacement requires revisiting that proof.\par
The special regime \(\sigma^2=\Theta(n)\) mentioned after Example 3.5 has \(\gamma=\Theta(1)\), and ordinary Gram concentration can suffice. It does not cover the permitted \(\sigma^2=1\) family or the entire \(\sigma^2=o(n)\) regime in Example 4.2.\par
The final risk formulas \(\mathrm{Err}_w\approx1-\eta\) and \(\mathrm{Err}_{w2s}\approx(1-\eta)^2\) may survive through a different or sharper analysis. They are not disproved here. The refuted assertion is decomposability under the example’s original full parameter range, which is the stated gateway to those conditional results.\par
\#\# Supported impact and limited repair\par
The supported impact is Major: a demonstrably invalid applicability derivation is used to justify the advertised benign-overfitting application of the framework. Neither Theorem 3.8 under genuine Assumption 3.7 nor the measured empirical correlations are refuted by this argument.\par
Equation reference: eq:repair. For the displayed Gaussian family with \(\sigma^2=1\) and \(k=1\), the stronger dimension condition \(m/n^3\to\infty\) restores the failed condition. Applying the Gaussian covariance bound recorded as Lemma C.3 to the transposed matrix, now with \(m\) samples and dimension \(n\), gives with probability at least \(1-2e^{-n/2}\) \[\left\|Z^\top Z/m-I\right\|_{\mathrm{op}}
 \le4\sqrt{n/m}+4n/m.\] Hence \[
 \frac{\|A_n\|_{\mathrm{op}}}{\gamma^2}
 \le4\sqrt{\frac{n^3}{m}}+4\frac{n^2}{m}\longrightarrow0.\] The other conditions checked above remain valid under this stronger scaling. This verifies restored decomposability for this specific Gaussian family; it does not verify a repair of the full sub-Gaussian theorem, arbitrary label distributions, or the unrestricted original parameter range.\par
Problem 2 (The finite inverse-perturbation bound omits a denominator).  Lemma A.3, PDF p. 13, states \(\|(A+\Delta)^{-1}-A^{-1}\|_{\mathrm{op}}\le
\|A^{-1}\|_{\mathrm{op}}^2\|\Delta\|_{\mathrm{op}}\) for invertible \(A\). This finite inequality is false.\par
\#\# Complete scalar check and correction\par
Equation reference: eq:resolvent. Let \(A=1\) and \(\Delta=-1/2\). Both \(A\) and \(A+\Delta\) are invertible. The stated right side equals \(1/2\), but \[|(1-1/2)^{-1}-1|=1.\] The resolvent identity instead yields \[(A+\Delta)^{-1}-A^{-1}
 =-(I+A^{-1}\Delta)^{-1}A^{-1}\Delta A^{-1}.\] If \(\|A^{-1}\Delta\|_{\mathrm{op}}<1\), the Neumann series gives \[
 \|(A+\Delta)^{-1}-A^{-1}\|_{\mathrm{op}}
 \le\frac{\|A^{-1}\|_{\mathrm{op}}^2\|\Delta\|_{\mathrm{op}}}
 {1-\|A^{-1}\Delta\|_{\mathrm{op}}}.\] The scalar example attains this corrected bound.\par
\#\# Dependency and severity\par
Lemmas A.5 and A.6 invoke the false bound. Under their genuine hypotheses, however, \(\|A^{-1}\Delta\|_{\mathrm{op}}=o(1)\): perturbations are \(o((\gamma+\delta)^2)\) or \(o((\gamma+\delta+\rho)^2)\), while the relevant inverse scale is at most reciprocal to the corresponding sum up to a constant; those sums are bounded and \(\beta\ge\delta\). Thus the denominator in [eq:resolvent] (\#eq:resolvent) is \(1-o(1)\), preserving the required absolute \(o(1)\) inverse errors. This supported Minor finite-bound defect therefore does not independently invalidate the inspected asymptotic conclusions. The grading follows this verified consequence.\par
\# Empirical interpretation and remaining boundaries\par
Corollaries 5.1–5.2 give operator-norm upper bounds, not universal equalities between a label-free score and prediction gap. Their inspected inequalities follow from submultiplicativity and the residual term under the stated assumptions. The paper separately tests finite nonlinear correlations.\par
Appendix D.2, PDF p. 37, discloses dataset-specific hyperparameter selection and best-result reporting in experiment III, as well as labeled performance measurements for tuning. Table 2 reports held-out-model tuning transfer for Justice and Commonsense in experiment II: \(0.885\pm0.16\) and \(0.67\pm0.20\). This is counterevidence to a blanket claim that all correlations are unvalidated selection artifacts. The metric is label-free after calibration; the supplied experiment III results do not by themselves establish a fully untuned, label-free deployment procedure. Raw metric/error pairs were unavailable for recomputation.\par
Theorem 3.6’s exact eigenspace preservation remains an additional unresolved boundary: Appendix C.3 bounds, but does not set to zero, \(\mathbb E[h\xi^\top]\), which does not by itself establish that the proposed subspace is a population eigenspace. No separate severity is assigned because intended dependence assumptions and the sub-Gaussian parameter convention need clarification. This does not increase Problem 1’s supported impact.\par
\# Conclusion and recorded state\par
Dimension | Recorded state\par
Inspection outcome | Established Gaussian applicability and finite perturbation-bound defects.\par
Supported impact | Major for Problem 1; Minor for Problem 2.\par
Repair verification | Stronger Gaussian scaling verified for the displayed family; inverse-bound denominator verified.\par
Author adoption | Unknown; no contact or adoption request.\par
Scientific adjudication | Pending; no human review supplied.\par
Remaining work | General-range repair, unresolved construction assumptions, empirical reproduction and later-version check.\par
1 Y. Xue, J. Li, and B. Mirzasoleiman, Representations Shape Weak-to-Strong Generalization: Theoretical Insights and Empirical Predictions, Proceedings of Machine Learning Research 267 (2025), 70120–70159. Reviewed published proceedings PDF, 40 pages. https://proceedings.mlr.press/v267/xue25d.html.\par
\# Independent computation and evidence binding\par
Problem 1’s chi-squared argument is analytic. The independent script uses NumPy 2.4.3 on Python 3.13.5, seed 20261024, float64 arithmetic. It constructs Gaussian matrices only for \(n=8,16,32\). Diagnostics for \(n=100,1000,10000\) instead sample the exact chi-squared representation and do not allocate \(n^3\)-row matrices. Actual results and script/source hashes are in . No author code, datasets, weights, or paid services were executed or obtained.\par
The actual original PDF SHA-256 is\par
> .\par
Source renders for PDF pp. 4, 5, 32 and 34 are preserved in this paper’s verification directory. The critical squared rate on p. 4 and step after equation (34) on p. 32 were visually inspected by the auxiliary reviewer. Parent integration, final evidence checking, final compilation and report-PDF visual inspection are recorded with their actual completion states in provenance rather than presumed here.\par
\# AI and human roles; scope ledger\par
The native auxiliary task `theory\_representations\_privacy` read the source, developed the arguments and independent diagnostic, and drafted the report and JSON records. No runtime/UI model identity was exposed to that task; the inherited instruction label was GPT-6 Astra Pro, not a separate observation of the visible parent model. Native auxiliary task subsequently performed an independent read-only challenge: it checked the Gaussian premises, exact population eigenspace, covariance boundedness, column lower bound, source dependency and scalar inverse correction, and visually read source pages 4 and 32. It did not rerun the simulation or audit every main-theorem proof. Its runtime identifier was likewise not exposed. The findings survived that challenge and the parent’s integration checks. No independent chat was opened by either auxiliary task and no human scientific review was supplied.\par
The coverage ledger separates inspected principal arguments from partially inspected general concentration proofs and unperformed empirical work. Bibliographic completeness, recomputation of plots, all possible defects, later versions and unknown OpenReview-only material remain outside the verified scope.\par
\end{document}
